\section{Festival and Alternative Launch}
This is an audience and programming hypothesis, not a researched calendar. No festival has been contacted and no selection, premiere status or eligibility is implied. A producer must verify current rules before submissions.

The film may fit programs for intimate international drama, work and social institutions, ensemble performance or regional independent filmmaking. A submission should foreground its distinct formal proposition: a forty-eight-hour practical task completed truthfully at an economic cost. ``An important story about archives'' is too vague.

If a festival route is unsuitable, a later curated theatrical or community screening path could pair the film with a moderated discussion about records access, labor and consent. Partners must be invited on their own terms and compensated where appropriate. The fictional town must not be presented as an allegation about a real authority.

Next decisions depend on the completed film, actual runtime, premiere commitments, distributor involvement and resources. Prepare captions, accessible screening information and a factual production description. Do not spend submission money on the strength of this generated strategy alone.
